\section{Setup and conventions}

% In this document we work with a fixed loopless matroid $M$ of rank $r$ on the
% ground set $[n]$, equipped with an order-reversing involution
% $F\mapsto F^\perp$ on its lattice of flats $\cLL(M)$. Our aim is to construct an
% order-reversing involution $\Phi$ on the relative Dressian $\Dr(M)$ extending
% the involution on $\cLL(M)^{\op}$.

\begin{convention}
Throughout, $M$ is a loopless matroid of rank $r$ on the ground set $E=[n]$,
with lattice of flats $\cLL(M)$ ordered by inclusion, rank function
$\rk_M$, and closure operator $\cl_M$. We assume $M$ is \emph{modular}, i.e.
there is an order-reversing involution
\[
\perp:\ \cLL(M)\longrightarrow \cLL(M),\qquad F\mapsto F^\perp,
\]
satisfying $F\subseteq G\iff G^\perp\subseteq F^\perp$, $(F^\perp)^\perp=F$, and
necessarily
\[
\rk_M(F)+\rk_M(F^\perp)=r,\qquad \hat 0^\perp=\hat 1,\ \hat 1^\perp=\hat 0 .
\]
We write $\Rbar=\R\cup\{+\infty\}$ with the tropical (min,+) conventions. All
valuated matroids below are quotients of the trivial valuation $\underline M$ of
$M$, i.e. they are supported on subsets of bases-extensions of $M$.
\end{convention}

% We first recall the precise definitions of valuated matroid quotients and of
% the relative Dressian, then give a cryptomorphic description of quotients in
% terms of flats, and finally use the involution on flats to define $\Phi$.

\begin{definition}\label{def:quotient}
Let $\theta_2$ be a valuated matroid of rank $r$ on $[n]$ and $\theta_1$ a
valuated matroid of rank $k\le r$ on $[n]$. We say $\theta_1$ is a
\emph{quotient} of $\theta_2$, written $\theta_2\onto\theta_1$, if for every
$S\subset[n]$ with $|S|=k-1$ and every $T\subset[n]$ with $|T|=k+1$ the tropical
\emph{incidence--Pl\"ucker relation}
\begin{equation}\label{eqn:incidence-plucker}
\min_{a\in T\smallsetminus S}\bigl\{\theta_1(S+a)+\theta_2^{*}(T-a)\bigr\}
\end{equation}
vanishes tropically, where $\theta_2^{*}$ is the valuated dual of $\theta_2$ of
rank $n-r$, evaluated on the appropriate complementary set; equivalently (and
this is the form we use) for all such $S,T$,
\begin{equation}\label{eqn:ip2}
\min_{a\in T\smallsetminus S}\bigl\{\theta_1(S+a)+\theta_2(\,\overline{T-a}\,)\bigr\}
\quad\text{attains its minimum at least twice or is $+\infty$,}
\end{equation}
with $\overline{(\cdot)}$ denoting the standard complementation for the dual.
\end{definition}

% For the trivial valuation $\underline M$ of $M$, the support condition imposed
% by \eqref{eqn:ip2} is exactly the matroid-quotient condition between the
% underlying matroids; thus a rank-$k$ quotient of $\underline M$ is a valuated
% matroid $\theta$ of rank $k$ whose underlying matroid is a quotient of $M$.

\begin{definition}
For a matroid $M$ as above and $0\le k\le r$, let $\bV^k(M)$ be the set of
rank-$k$ valuated matroid quotients of $\underline M$, and let
$\bV(M)=\bigsqcup_{k=0}^{r}\bV^k(M)$, partially ordered by the quotient relation
$\onto$. After the standard tropical rescaling, $\bV(M)/\R\cong\Dr(M)$ as posets
(quotient $\onto$ corresponds to inclusion $\supseteq$ of tropical linear
spaces).
\end{definition}

\section{A cryptomorphism: quotients as functions on flats}

% The key technical input is that a valuated matroid quotient $\theta$ of
% $\underline M$ is determined by, and freely encoded by, its restriction to the
% rank-$k$ flats of $M$, subject to a system of "essential" tropical relations
% indexed by modular pairs of flats. This is the matroid analogue of the
% classical statement that the Pl\"ucker coordinates of a coordinate subspace are
% indexed by flats.

\begin{lemma}\label{lem:flat-support}
Let $\theta\in\bV^k(M)$ and let $N$ be the underlying rank-$k$ matroid of
$\theta$, a quotient of $M$. For a $k$-subset $B\subseteq[n]$ we have
$\theta(B)<+\infty$ only if $\cl_M(B)$ has $M$-rank $\ge k$, and for two
$k$-subsets $B,B'$ with $\cl_M(B)=\cl_M(B')=:F$ of $M$-rank exactly $k$ and both
bases of $N$, the differences $\theta(B)-\theta(B')$ are determined by the local
data of $\theta$ at $F$.
\end{lemma}

\begin{proof}
Since $N$ is a quotient of $M$, every flat of $N$ is a flat of $M$, and
$\rk_N\le \rk_M$. If $B$ is a basis of $N$ then $\cl_M(B)\supseteq \cl_N(B)=[n]$
is impossible unless $k=r$; in general $B$ independent in $N$ forces $B$
independent in $M$ (as $N$ is a quotient and quotients only increase dependency
in the reverse direction; more precisely independence in the quotient implies
independence upstairs). Hence $\cl_M(B)$ has $M$-rank $\ge k$, proving the first
claim.

For the second claim, fix a flat $F$ of $M$ with $\rk_M(F)=k$. Consider two
$k$-subsets $B,B'\subseteq F$ that are bases of $N$ with $\cl_M(B)=\cl_M(B')=F$.
By the basis-exchange chain in the matroid $M|_F$ (which has rank $k$ and is
connected enough to exchange elements one at a time within $F$), we may pass
from $B$ to $B'$ by a sequence of single-element exchanges
$B=B_0,B_1,\dots,B_m=B'$ with $|B_i\triangle B_{i+1}|=2$ and all $B_i\subseteq F$
of rank $k$ in $M$. For each such elementary exchange, writing
$B_i=S+a$, $B_{i+1}=S+b$ with $|S|=k-1$, the tropical Pl\"ucker relation for
$\theta$ labeled by $(S, S\cup\{a,b\}\cup\{c\})$ — using a third element $c$ of
the universal extension or, when $k=r$, the relation interior to $F$ — shows the
value $\theta(S+a)-\theta(S+b)$ is a difference of finite quantities determined
by $\theta$ restricted to $k$-subsets of $F$. Iterating, all values $\theta(B)$
for $B$ a basis of $N$ with $\cl_M(B)=F$ are determined, up to a single additive
constant $c_F\in\R$, by $\theta|_F$. This is exactly the assertion.
\end{proof}

\begin{definition}\label{def:flat-function}
For $\theta\in\bV^k(M)$ define a function
\[
\widehat\theta:\ \cLL_k(M)\longrightarrow\Rbar,\qquad
\widehat\theta(F):=\min_{\substack{B\subseteq F,\ |B|=k\\ \cl_M(B)=F}}\theta(B),
\]
on the set $\cLL_k(M)$ of $M$-flats of rank $k$ (and $\widehat\theta(F)=+\infty$
if no $k$-subset $B\subseteq F$ with $\theta(B)<\infty$ exists).
\end{definition}

% The next proposition is the cryptomorphism alluded to in the problem
% statement. It identifies $\bV^k(M)$ with the set of functions on rank-$k$ flats
% satisfying a system of "essential" tropical Pl\"ucker relations indexed by
% covering modular pairs.

\begin{proposition}[Essential relations]\label{prop:essential-relations}
The assignment $\theta\mapsto\widehat\theta$ is a bijection (after tropical
rescaling) between $\bV^k(M)$ and the set $\bP^k(M)$ of functions
$p:\cLL_k(M)\to\Rbar$, not identically $+\infty$, satisfying the
\emph{essential relations}: for every $G\in\cLL_{k-1}(M)$ and every
$H\in\cLL_{k+1}(M)$ with $G\subset H$, the expression
\begin{equation}\label{eqn:essential}
\min_{\substack{F\in\cLL_k(M)\\ G\subset F\subset H}} p(F)
\end{equation}
attains its minimum at least twice (or is $+\infty$). Moreover, $\theta_2\onto
\theta_1$ in $\bV(M)$ if and only if for all $G\in\cLL_{k-1}(M)$ (with
$k=\rk\theta_1$) and all $H\in\cLL_{r-?}\dots$ the corresponding incidence
relations among $\widehat\theta_1,\widehat\theta_2$ hold; concretely, the order
on $\bV(M)$ is reflected faithfully by the essential incidence relations on the
flat-functions.
\end{proposition}

\begin{proof}
We prove that $\theta\mapsto\widehat\theta$ is well defined, injective, and
surjective onto $\bP^k(M)$.

\emph{Well-defined and the essential relations hold.}
Let $\theta\in\bV^k(M)$. Fix $G\in\cLL_{k-1}(M)$ and $H\in\cLL_{k+1}(M)$ with
$G\subset H$, so $\rk_M(H)-\rk_M(G)=2$. Choose a $(k-1)$-subset
$S\subseteq G$ with $\cl_M(S)=G$, and choose two elements
$a,b\in H\smallsetminus G$ with $\cl_M(G\cup\{a,b\})=H$. The tropical Pl\"ucker
relation for $\theta$ labeled by $(S,\,S\cup\{a,b\})$ — together with the support
restriction from \Cref{lem:flat-support}, which forces all finite contributions
to come from $k$-subsets whose $M$-closure is a rank-$k$ flat $F$ with
$G\subset F\subset H$ — yields that
\[
\min_{a'\in (S\cup\{a,b\})\smallsetminus S}\{\theta(S+a')\}
=\min_{\substack{F\in\cLL_k(M)\\ G\subset F\subset H}}\widehat\theta(F)
\]
attains its minimum at least twice, since by modularity of the interval
$[G,H]$ in $\cLL(M)$ (an interval of length two in any geometric lattice is a
modular lattice, in fact a line) the rank-$k$ flats $F$ with $G\subset F\subset
H$ are exactly the atoms of the rank-two interval $[G,H]$, and there are at least
two of them whenever any term is finite. Independence of the chosen $S$ and of
the chosen representatives follows from \Cref{lem:flat-support}: changing
representatives shifts all finite values by a common additive constant, which
does not affect the "minimum attained twice" condition. Hence
$\widehat\theta\in\bP^k(M)$.

\emph{Injectivity.}
By \Cref{lem:flat-support}, for each rank-$k$ flat $F$ the values
$\{\theta(B): \cl_M(B)=F,\ B\text{ basis of }N\}$ are determined up to one
additive constant $c_F$, and $\widehat\theta(F)$ records this constant
(the minimum over $B\subseteq F$). The local pattern within $F$ (i.e. the
$\theta(B)$ for fixed $F$) is itself a valuated matroid on $M|_F$, namely the
trivial valuation of the matroid quotient $N|_F$; this local datum is determined
by $N$, which in turn is determined by the support of $\widehat\theta$ via the
matroid cryptomorphism: the rank-$k$ flats $F$ with $\widehat\theta(F)<\infty$
are exactly the flats of $N$ of rank $k$, and a matroid is determined by its
flats. Therefore $\theta$ is recovered from $\widehat\theta$ up to the global
tropical rescaling, proving injectivity in $\bV^k(M)/\R$.

\emph{Surjectivity.}
Let $p\in\bP^k(M)$. Define $N$ to be the matroid whose rank-$k$ flats are the
$F$ with $p(F)<\infty$ together with the closures forced by the lattice
structure; the essential relations \eqref{eqn:essential} guarantee that this
collection is closed under the geometric-lattice axioms (every length-two
interval has at least two atoms whenever it has one finite atom, which is exactly
the exchange axiom for flats), so $N$ is a well-defined matroid quotient of $M$.
Now set, for each $k$-subset $B$ with $\cl_M(B)=F$ a basis of $N$,
\[
\theta(B):=p(F)+\lambda_{N|_F}(B),
\]
where $\lambda_{N|_F}$ is the (unique up to constant) local valuation coming from
the trivial valuation of $N|_F$, and $\theta(B):=+\infty$ otherwise. The
three-term Pl\"ucker relations for $\theta$ split into two types: those internal
to a fixed rank-$k$ flat $F$ (which hold because $\lambda_{N|_F}$ is a valuated
matroid) and those crossing a covering pair $(G,H)$ as above (which hold exactly
because $p$ satisfies the essential relation \eqref{eqn:essential}). Hence
$\theta\in\bV^k(M)$ and $\widehat\theta=p$.

\emph{Compatibility with the order.}
If $\theta_2\onto\theta_1$ with $\rk\theta_2=r$, $\rk\theta_1=k$, then the
incidence relations \eqref{eqn:ip2} restrict, by the same support analysis and
modularity of length-two intervals, to incidence relations between
$\widehat\theta_1$ and $\widehat\theta_2$ on flats; conversely these incidence
relations on flats reconstruct \eqref{eqn:ip2} by the local-valuation
factorization above. Thus the quotient order on $\bV(M)$ is faithfully encoded by
the essential incidence relations on flat-functions.
\end{proof}

\section{The duality involution}

% With the cryptomorphism in hand, the order-reversing involution on flats
% transports to an involution on flat-functions, hence on $\bV(M)$ and on
% $\Dr(M)$.

\begin{definition}\label{def:Phi}
For $\theta\in\bV^k(M)$ with flat-function $\widehat\theta\in\bP^k(M)$, define
$\theta':=\Phi(\theta)$ to be the unique element of $\bV^{r-k}(M)$ whose
flat-function is
\[
\widehat{\theta'}:\ \cLL_{r-k}(M)\longrightarrow\Rbar,\qquad
\widehat{\theta'}(F):=\widehat\theta(F^\perp),
\]
provided the right-hand side lies in $\bP^{r-k}(M)$ (this is verified in
\Cref{thm:main}). Equivalently, in Pl\"ucker coordinates $\Phi$ is the tropical
linear map $\theta\mapsto\theta'$ with $\theta'(F)=\theta(F^\perp)$.
\end{definition}

\begin{lemma}\label{lem:perp-preserves-relations}
The reindexing $p\mapsto p^\perp$, $p^\perp(F):=p(F^\perp)$, is a bijection
$\bP^k(M)\to\bP^{r-k}(M)$.
\end{lemma}

\begin{proof}
Since $\perp$ is a bijection $\cLL_k(M)\to\cLL_{r-k}(M)$ (it is rank-additive:
$\rk_M(F)+\rk_M(F^\perp)=r$), the map $p\mapsto p^\perp$ is a bijection between
functions on $\cLL_k(M)$ and functions on $\cLL_{r-k}(M)$. It remains to verify
the essential relations are preserved.

Fix $G'\in\cLL_{(r-k)-1}(M)$ and $H'\in\cLL_{(r-k)+1}(M)$ with $G'\subset H'$;
we must show
\[
\min_{\substack{F'\in\cLL_{r-k}(M)\\ G'\subset F'\subset H'}} p^\perp(F')
=\min_{\substack{F'\in\cLL_{r-k}(M)\\ G'\subset F'\subset H'}} p(F'^{\perp})
\]
attains its minimum at least twice (or is $+\infty$). Apply $\perp$: as $\perp$
is an order-reversing bijection on $\cLL(M)$, the correspondence $F'\mapsto
F'^\perp$ carries the set
\[
\{F'\in\cLL_{r-k}(M): G'\subset F'\subset H'\}
\]
bijectively onto
\[
\{F\in\cLL_{k}(M): H'^{\perp}\subset F\subset G'^{\perp}\}.
\]
Now $H'^{\perp}\in\cLL_{k-1}(M)$, $G'^{\perp}\in\cLL_{k+1}(M)$, and
$H'^{\perp}\subset G'^{\perp}$ is again a length-two interval (since
$\rk_M(G'^\perp)-\rk_M(H'^\perp)= (k+1)-(k-1)=2$). Under this bijection the
multiset of values
\[
\{p^\perp(F'): G'\subset F'\subset H'\}=\{p(F): H'^\perp\subset F\subset G'^\perp\},
\]
so the minimum on the left attains its minimum at least twice if and only if the
minimum on the right does. But the right-hand side is precisely the essential
relation for $p\in\bP^k(M)$ at the covering pair
$(G,H)=(H'^{\perp},G'^{\perp})$, which holds by hypothesis. Hence $p^\perp$
satisfies all essential relations and $p^\perp\in\bP^{r-k}(M)$.

The map $p\mapsto p^\perp$ is involutive because $(F^\perp)^\perp=F$, so it is a
bijection $\bP^k(M)\leftrightarrow\bP^{r-k}(M)$.
\end{proof}

\begin{theorem}\label{thm:main}
Let $M$ be a modular matroid of rank $r$ with order-reversing involution
$F\mapsto F^\perp$ on $\cLL(M)$. Then the map $\Phi:\bV(M)\to\bV(M)$ of
\Cref{def:Phi} is a well-defined, order-reversing involution that descends to an
order-reversing involution
\[
\Phi:\ \Dr(M)\longrightarrow\Dr(M)
\]
on the relative Dressian, and $\Phi$ extends the involution
$F\mapsto F^\perp$ under the natural inclusion $\cLL(M)^{\op}\subset\Dr(M)$.
\end{theorem}

\begin{proof}
\emph{Well-definedness and involutivity.}
By \Cref{prop:essential-relations}, $\theta\mapsto\widehat\theta$ identifies
$\bV^k(M)$ (mod rescaling) with $\bP^k(M)$. By
\Cref{lem:perp-preserves-relations}, $p\mapsto p^\perp$ is a bijection
$\bP^k(M)\to\bP^{r-k}(M)$ that is an involution on
$\bigsqcup_k\bP^k(M)$. Transporting through the cryptomorphism gives a well-defined
involution $\Phi$ on $\bV(M)$ with $\widehat{\Phi(\theta)}=(\widehat\theta)^\perp$.
Since $((\widehat\theta)^\perp)^\perp=\widehat\theta$ by
$(F^\perp)^\perp=F$, we have $\Phi\circ\Phi=\mathrm{id}$.

\emph{Order-reversal.}
Suppose $\theta_2\onto\theta_1$ in $\bV(M)$, with $\rk\theta_2=k_2\ge
k_1=\rk\theta_1$. By the compatibility statement in
\Cref{prop:essential-relations}, this is equivalent to a system of incidence
relations between $\widehat\theta_2$ and $\widehat\theta_1$ indexed by
covering pairs $(G,H)$ with $G\in\cLL_{k_1-1}(M)$ and the relevant flats. Applying
the bijection $\perp$ to all indices, exactly as in the proof of
\Cref{lem:perp-preserves-relations}, transforms this system into the incidence
system characterizing $\Phi(\theta_1)\onto\Phi(\theta_2)$: indeed $\perp$ sends a
rank-$k_1$ flat to a rank-$(r-k_1)$ flat and reverses all inclusions, so the
incidence relations witnessing "$\theta_1$ is a quotient encoded below
$\theta_2$" become precisely those witnessing "$\Phi(\theta_2)$ is a quotient of
$\Phi(\theta_1)$." Concretely, since $\widehat{\Phi(\theta_i)}(F)=
\widehat{\theta_i}(F^\perp)$ and $\perp$ is an order-reversing bijection of
$\cLL(M)$ preserving the multisets of values appearing in each essential
incidence relation, every incidence relation valid for the pair
$(\theta_2,\theta_1)$ is valid, after reindexing by $\perp$, for the pair
$(\Phi(\theta_1),\Phi(\theta_2))$. Hence
\[
\theta_2\onto\theta_1\quad\Longrightarrow\quad \Phi(\theta_1)\onto\Phi(\theta_2).
\]
Because $\Phi$ is an involution, the converse implication also holds, so
$\Phi$ is order-reversing. Since the partial order on $\Dr(M)$ is the quotient
order $\onto$ read as inclusion $\supseteq$ of tropical linear spaces, $\Phi$
descends to an order-reversing involution on $\Dr(M)$ (it is constant on
tropical-rescaling classes because $p\mapsto p^\perp$ commutes with adding a
global constant).

\emph{Extension of $F\mapsto F^\perp$.}
Recall the inclusion $\cLL(M)^{\op}\subset\Dr(M)$ sends a flat $F$ to
$\Trop(M/F\oplus U_{0,F})$, the tropical linear space of the trivial valuation of
the matroid quotient $M/F\oplus U_{0,F}$ of rank $r-\rk_M(F)$; call the
corresponding element $\theta_F\in\bV^{r-\rk_M(F)}(M)$. Its flat-function is, by
construction,
\[
\widehat{\theta_F}(K)=
\begin{cases}
0,& F\subseteq K,\ \rk_M(K)=r-\rk_M(F)\ \text{and } K \text{ is a flat of } M/F\oplus U_{0,F},\\
+\infty,& \text{otherwise,}
\end{cases}
\]
i.e. $\widehat{\theta_F}$ is the $0/\infty$ indicator of the rank-$(r-\rk_M F)$
flats lying above $F$. Now compute $\Phi(\theta_F)$: its flat-function on a
rank-$\rk_M(F)$ flat $K$ is
\[
\widehat{\Phi(\theta_F)}(K)=\widehat{\theta_F}(K^\perp).
\]
This equals $0$ exactly when $F\subseteq K^\perp$, i.e. (applying $\perp$ and
using order-reversal) exactly when $K\subseteq F^\perp$, and $+\infty$
otherwise. Therefore $\widehat{\Phi(\theta_F)}$ is the $0/\infty$ indicator of
the rank-$\rk_M(F)=r-\rk_M(F^\perp)$ flats contained in $F^\perp$, which is
precisely the flat-function $\widehat{\theta_{F^\perp}}$ of the element
$\theta_{F^\perp}=\Trop(M/F^\perp\oplus U_{0,F^\perp})$. By the bijectivity of
$\theta\mapsto\widehat\theta$ (\Cref{prop:essential-relations}),
\[
\Phi(\theta_F)=\theta_{F^\perp}.
\]
Thus under the inclusion $\cLL(M)^{\op}\subset\Dr(M)$ the involution $\Phi$
restricts to $F\mapsto F^\perp$, as required.
\end{proof}

\begin{corollary}
For any modular matroid $M$ with chosen order-reversing involution on
$\cLL(M)$, there exists an order-reversing involution $\Phi$ on the relative
Dressian $\Dr(M)$ extending $F\mapsto F^\perp$. In particular this answers the
Question affirmatively.
\end{corollary}

\begin{proof}
Immediate from \Cref{thm:main}.
\end{proof}